% =============================================================
%  CV — ZINEB MEFTAH  (colonne unique)
%  Compiler avec : tectonic CV_Zineb_Meftah_FR.tex
%
%  POURQUOI UNE SEULE COLONNE (2026-09-19) — la version deux colonnes est conservée
%  dans documents/_legacy_two_column/. Un ATS lit le CV LINÉAIREMENT, et la barre
%  latérale s'intercalait au milieu du contenu principal : « EXPÉRIENCE
%  PROFESSIONNELLE » était immédiatement suivi de « LANGUES », le titre du poste
%  était séparé de l'employeur, et une phrase était coupée en deux par le tableau
%  des langues (« ... en langage » / six lignes de niveaux / « naturel de
%  documentations ... »). Aucun analyseur ne pouvait rattacher GE HealthCare au
%  poste. En colonne unique, l'ordre de lecture du PDF EST l'ordre visuel.
% =============================================================
\documentclass[10pt, a4paper]{article}

\usepackage[T1]{fontenc}
\usepackage{microtype}
% POLICES OPENTYPE, ET SANS LES LIGATURES f. Ce sont les mêmes caractères que mathpazo et
% tgheros dessinaient (Pagella = Palatino, Heros = Helvetica), mais chargés en OpenType, seule
% façon de désactiver les ligatures sous XeTeX — `\DisableLigatures` de microtype est réservé à
% pdfTeX, et tectonic tourne sur XeTeX.
%
% POURQUOI ÇA COMPTE : une ligature entre dans la COUCHE TEXTE du PDF, pas seulement dans le
% dessin. Mesuré le 2026-09-23 sur un CV réellement envoyé : « Intelligence Artificielle » en
% ressortait « Intelligence Artificielle » (U+FB01), et de même « certification »,
% « qualification », « classification », « fine-tuning ». Un ATS qui cherche littéralement
% « Artificielle » — le mot-clé central de ses candidatures IA — ne trouvait rien du tout.
% Les fichiers sont désignés par leur NOM DE FICHIER : le nom de famille (« TeX Gyre Pagella »)
% n'est pas résolu par tectonic, qui sert les polices depuis son propre paquetage.
\usepackage{fontspec}
\setmainfont{texgyrepagella-regular.otf}[BoldFont=texgyrepagella-bold.otf,%
  ItalicFont=texgyrepagella-italic.otf,Ligatures=NoCommon]
\setsansfont{texgyreheros-regular.otf}[BoldFont=texgyreheros-bold.otf,Ligatures=NoCommon]
\renewcommand{\familydefault}{\rmdefault}
\newcommand{\headfont}{\sffamily}
\usepackage{fontawesome5}
\usepackage{geometry}
\usepackage{xcolor}
\usepackage{array}
\usepackage{tabularx}
\usepackage{enumitem}
\usepackage{hyperref}
\usepackage{parskip}
\usepackage{calc}

\geometry{top=11mm,bottom=9mm,left=16mm,right=16mm,nohead,nofoot}
\pagestyle{empty}
\setlength{\parindent}{0pt}
\setlength{\parskip}{0pt}
\setlength{\topskip}{0pt}

\definecolor{accent}{HTML}{1E3A8A}
\definecolor{subtitleColor}{HTML}{6B7280}
\definecolor{gold}{HTML}{9CA3AF}
\definecolor{headingColor}{HTML}{1E3A8A}
\definecolor{ruleColor}{HTML}{D1D5DB}
\definecolor{mutedText}{HTML}{6B7280}
% Titres de section : volontairement PLUS LÉGERS que les mots-clés du contenu. L'œil doit tomber
% sur « Deep Learning », pas sur le mot « FORMATION », qui n'apprend rien à personne.
\definecolor{sectionColor}{HTML}{1E3A8A}
\definecolor{linkColor}{HTML}{1E3A8A}

\hypersetup{colorlinks=true,urlcolor=linkColor,linkcolor=linkColor,pdfborder={0 0 0}}

\newlength{\mainInner}\setlength{\mainInner}{\textwidth}

% --- AUTO-FIT ---------------------------------------------------------------
% Identique à la version deux colonnes : le contenu est mesuré par TeX lui-même et
% cv_builder recompile à \cvFit plus serré, puis retire des blocs, jusqu'à ce qu'il
% tienne. pdftotext ne peut pas faire ce travail — le texte dessiné sous le bord de
% la page disparaît de sa sortie, donc il sous-estime le débordement.
\providecommand{\cvFit}{1.0}
\newsavebox{\cvMainBox}
\newlength{\cvMainHeight}
\newlength{\cvMainAvail}
\setlength{\cvMainAvail}{\dimexpr\paperheight-11mm-9mm\relax}
\newlength{\cvgap}
\newcommand{\vgap}[1]{\setlength{\cvgap}{#1}\vspace{\cvFit\cvgap}}

% ⚠ L'INTERLETTRAGE EST INTERDIT ICI, et c'est une leçon payée en le mesurant (2026-10-02).
% `\addfontfeature{LetterSpace=12}` donne bien un titre plus léger à l'œil, mais la couche texte
% du PDF rend alors « FO R M AT I O N » : un ATS ne retrouve plus un seul mot du titre. La même
% famille de bug que la fusion icône+valeur réglée par \ic. La hiérarchie passe donc par la
% TAILLE et la COULEUR, qui ne touchent pas à l'extraction : titre 10pt gris ardoise contre
% mots-clés 8,8pt noirs en gras — le gras du contenu domine, le titre recule.
\newcommand{\mainSection}[1]{%
  \vgap{6pt}%
  {\headfont\color{sectionColor}\fontsize{10pt}{11.5pt}\selectfont\bfseries\MakeUppercase{#1}}\par%
  \nointerlineskip\vgap{2pt}%
  {\color{ruleColor}\rule{\mainInner}{0.8pt}}\par\vgap{4pt}%
}

\newcommand{\cvEntry}[4]{%
  \begin{tabularx}{\mainInner}{@{}X r@{}}
    {\headfont\fontsize{10pt}{11.5pt}\selectfont #1} &
    {\headfont\color{mutedText}\fontsize{8.5pt}{10pt}\selectfont #3}
  \end{tabularx}\par\vspace{1pt}%
  {\color{subtitleColor}\fontsize{9.5pt}{10pt}\selectfont\itshape #2}\par
  #4\vgap{6pt}%
}

\newcommand{\cvOneLine}[2]{%
  \begin{tabularx}{\mainInner}{@{}X r@{}}
    {\headfont\fontsize{10pt}{11.5pt}\selectfont #1} &
    {\headfont\color{mutedText}\fontsize{8.5pt}{10pt}\selectfont #2}
  \end{tabularx}\par\vgap{4pt}%
}

\newcommand{\cvProject}[3]{%
  \begin{tabularx}{\mainInner}{@{}X r@{}}
    {\headfont\fontsize{10pt}{11.5pt}\selectfont #1} &
    {\headfont\color{mutedText}\fontsize{8.5pt}{10pt}\selectfont #2}
  \end{tabularx}\par\vspace{1pt}%
  #3\vgap{5pt}%
}

\newcommand{\cvBullets}[1]{%
  \vspace{2pt}%
  \begin{itemize}[leftmargin=10pt,itemsep=1pt,topsep=0pt,partopsep=0pt,parsep=0pt,
    label={\color{gold}\small$\bullet$}]
    #1
  \end{itemize}%
}
\newcommand{\cvItem}[1]{\item{\fontsize{8.8pt}{10.8pt}\selectfont #1}}

% \ic : icône dorée dans une boîte fixe, SUIVIE D'UNE VRAIE ESPACE.
% L'espace n'est pas cosmétique : sans elle le glyphe et la valeur fusionnent en un
% seul token dans la couche texte du PDF, qui est ce que lit un ATS.
\newcommand{\ic}[1]{\makebox[4.2mm][l]{\color{gold}#1}\ }


% ══ AUCUNE CÉSURE, JAMAIS (2026-10-05) ══════════════════════════════════════════════════════
% Un mot coupé en fin de ligne devient « Transform-ers » dans la couche texte du PDF. pdftotext
% le recolle, mais tous les analyseurs d'ATS ne le font pas — et le mot coupé était justement
% « Hugging Face Transformers », c'est-à-dire exactement le genre de terme sur lequel une
% candidature est mise en correspondance. Une pénalité infinie interdit la coupure ; `tolerance`
% laisse LaTeX étirer les espaces à la place, ce qui est le bon compromis sur un CV (lignes
% courtes, vocabulaire technique) et non sur un livre.
\hyphenpenalty=10000
\exhyphenpenalty=10000
\tolerance=3000
\emergencystretch=3em

\begin{document}

\begin{lrbox}{\cvMainBox}%
\begin{minipage}[t]{\mainInner}

  {\headfont\fontsize{27pt}{30pt}\selectfont\bfseries ZINEB MEFTAH}\par\vspace{5pt}
  {\headfont\color{accent}\fontsize{12.5pt}{16pt}\selectfont\bfseries
    AI ENGINEER {\color{mutedText}$\cdot$} MLOPS {\color{mutedText}$\cdot$} DEEP LEARNING}\par\vspace{3pt}
  {\color{subtitleColor}\fontsize{9.5pt}{13pt}\selectfont\itshape
    Systèmes IA autonomes en production $\cdot$ CDI, CDD ou alternance\\
    Disponible dès oct. 2026 $\cdot$ Île-de-France ou télétravail}\par\vspace{5pt}

  % CONTACT — une seule ligne logique, juste sous l'en-tête, là où tout analyseur
  % de CV cherche les coordonnées. Chaque valeur est un token propre.
  {\fontsize{9pt}{13pt}\selectfont
    \ic{\faEnvelope}\href{mailto:you@example.com}{you@example.com}\hfill
    \ic{\faPhone}+33 6 00 00 00 00\hfill
    \ic{\faMapMarker}Île-de-France\hfill
    % PERMIS B, à sa demande (2026-10-01 : « c'est important »). Beaucoup d'annonces françaises
    % en font une condition explicite, et un recruteur qui ne le voit pas doit poser la question.
    % Posé ICI, sur la ligne de contact, pour deux raisons : les \hfill redistribuent, donc cela
    % ne coûte AUCUNE ligne — or une ligne de plus fait tomber un projet entier via l'auto-fit
    % (le rythme d'alternance avait fait passer le CV 1pt au-dessus et supprimé LeRobot) — et
    % c'est l'endroit où tout analyseur de CV cherche les informations personnelles.
    \ic{\faCar}Permis B\par\vspace{2pt}
    \ic{\faLinkedin}\textbf{\href{https://www.linkedin.com/in/zinebmeftah}{linkedin.com/in/zinebmeftah}}\hfill
    \ic{\faGithub}\href{https://github.com/ZinebMEFTAH}{github.com/ZinebMEFTAH}\hfill
    \ic{\faRobot}\href{https://huggingface.co/zino36}{huggingface.co/zino36}\hfill
    \ic{\faGlobe}\href{https://zinebmeftah.github.io}{zinebmeftah.github.io}\par
  }
  \vgap{4pt}

  \mainSection{Profil}
  {\fontsize{9pt}{12.5pt}\selectfont
% @profil
En \textbf{M1 MLSD} (Machine Learning pour la Science des Données) à l'Université Paris Cité. \textbf{Major de promotion} en licence IA (\textbf{1\textsuperscript{ère}/126}), après une classe préparatoire à \textbf{l'ENSIA}. Je conçois et j'exploite des \textbf{systèmes d'IA autonomes en production} : RAG, agents, MLOps. J'arbitre une architecture sur son coût et sa maintenance.
% @endprofil\par
  }
  \vspace{5pt}

  \mainSection{Compétences}
  {\fontsize{9pt}{12.5pt}\selectfont
% @skills
{\color{mutedText}Langages} \textbf{Python} (Expert), \textbf{C/C++}, \textbf{Java}, \textbf{JavaScript}, \textbf{SQL}\\
{\color{mutedText}LLM \& IA} \textbf{RAG}, \textbf{Agents}, \textbf{MCP}, \textbf{Claude API}, \textbf{Vector DB} $\cdot$ {\color{mutedText}ML / DL} \textbf{PyTorch}, \textbf{TensorFlow}, \textbf{Scikit-learn}, \textbf{Transformers}\\
{\color{mutedText}MLOps} \textbf{Docker}, \textbf{CI/CD}, \textbf{Git}, \textbf{Linux}, \textbf{Cloud} $\cdot$ {\color{mutedText}Données} \textbf{Pandas}, \textbf{NumPy}, \textbf{Embeddings}, \textbf{BM25}, \textbf{Reranking}\\
{\color{mutedText}Domaines} \textbf{NLP}, \textbf{Vision par ordinateur}, \textbf{Reinforcement Learning}, \textbf{Systèmes distribués} $\cdot$ {\color{mutedText}Web} \textbf{Flask}, \textbf{Node.js}, \textbf{Next.js}
% @endskills
\par
  }
  \vgap{2pt}

  \mainSection{Expérience Professionnelle}
  \cvEntry{GE HealthCare -- Buc (78), \^Ile-de-France}
    {Stagiaire Ingénieure IA \& MLOps}
    {Mai -- Juil. 2026}{%
% @bullets id=ge-healthcare-buc-78-ile-d
    \cvBullets{
      \cvItem{Agent \textbf{RAG} \textbf{validé par l'équipe et approuvé} pour la mise en production : interrogation en langage naturel de documentations techniques complexes, chaque réponse tracée jusqu'à la \textbf{section exacte} du document source, ramenant des recherches de plusieurs heures à quelques secondes. Pipeline codé de bout en bout en skills modulaires, puis ré-architecturé en \textbf{système multi-agents} sur \textbf{Microsoft Copilot Studio} quand le \textbf{coût en tokens} est devenu le facteur limitant}
    }
% @endbullets
  }

  \mainSection{Projets Sélectionnés}
% @cvblock id=outreach-agent rank=0 focus=ai,backend,mlops,data,fullstack keep=1
  \cvProject{Agent d'Outreach IA -- Système Autonome en Production}
    {2026 -- Présent}{%
% @bullets id=outreach-agent
    \cvBullets{
      \cvItem{Tourne \textbf{seul sur VM depuis mai 2026} : \textbf{3\,380 offres} suivies, \textbf{2\,150 entreprises} démarchées $\cdot$ vérifie chaque adresse, s'auto-limite contre le spam, apprend ce qui obtient des réponses $\cdot$ \textbf{Python} $\cdot$ \textbf{Claude API} $\cdot$ \textbf{Playwright} $\cdot$ \textbf{Google Cloud} $\to$ \href{https://github.com/ZinebMEFTAH/stationf-outreach-agent}{GitHub}}
    }
% @endbullets
  }
% @endcvblock

% @cvblock id=ams-monitoring rank=1 focus=mlops,backend
  \cvProject{AMS -- Monitoring Distribu\'e Multi-Serveurs}
    {Juin 2025}{%
% @bullets id=ams-monitoring
    \cvBullets{
      \cvItem{Supervision temps r\'eel de serveurs Linux : sondes CPU/RAM/disque/processus, d\'etection d'alertes, dashboard Flask $\cdot$ \textbf{9 t\^aches cron} (sondes, synchronisation distante, rotation des sauvegardes, purge des logs) $\cdot$ \textbf{Python} (psutil) $\cdot$ \textbf{Bash} $\cdot$ \textbf{Linux} $\cdot$ \textbf{SCP/SSH} $\to$ \href{https://github.com/ZinebMEFTAH/Servers-Monitoring-System}{GitHub}}
    }
% @endbullets
  }
% @endcvblock

% @cvblock id=content-engine rank=4 focus=fullstack,backend,data
  \cvProject{Moteur de Contenu Autonome -- Multi-Plateformes en Production}
    {2026 -- Présent}{%
% @bullets id=content-engine
    \cvBullets{
      \cvItem{Système en boucle fermée publiant sur \textbf{YouTube, TikTok, Instagram \& Facebook} : pilotage par email, scraping des vues/likes, republication hebdomadaire automatique des meilleurs titres $\cdot$ liens d'affiliation Amazon $\cdot$ \textbf{n8n} $\cdot$ \textbf{AWS EC2} $\cdot$ LLM}
    }
% @endbullets
  }
% @endcvblock

% @cvblock id=lerobot rank=3 focus=ai,mlops,data
  \cvProject{LeRobot -- Apprentissage Robotique \& MLOps}
    {Jan. -- Mars 2025}{%
% @bullets id=lerobot
    \cvBullets{
      \cvItem{Pipeline MLOps complet : entraînement d'agents robotiques (\textbf{Diffusion Models}, \textbf{ACT}) $\cdot$ \textbf{PyTorch} $\cdot$ \textbf{Hugging Face} $\to$ \href{https://huggingface.co/zino36/lerobot-pusht-colab}{Modèle} $\cdot$ \href{https://huggingface.co/spaces/zino36/lerobot-pusht-trainer}{Training UI}}
    }
% @endbullets
  }
% @endcvblock

% @cvblock id=compiler rank=2 focus=mlops,backend
  \cvProject{Compilateur Pascal $\to$ Assembleur x86-64}
    {Avril 2025}{%
% @bullets id=compiler
    \cvBullets{
      \cvItem{Cha\^ine de compilation \'ecrite de z\'ero : \textbf{Flex}, \textbf{Bison}, g\'en\'eration d'assembleur \textbf{x86-64} $\cdot$ \textbf{C++17} $\cdot$ \textbf{GNU/Linux} $\to$ \href{https://github.com/ZinebMEFTAH/MonCompilateur}{GitHub}}
    }
% @endbullets
  }
% @endcvblock

% @cvblock id=research-hf rank=2 focus=ai,data

  \cvProject{Recherche publiée -- ``Tags Generation Dataset'', Blog Hugging Face}
    {Août -- Déc. 2024}{%
% @bullets id=research-hf
    \cvBullets{
      \cvItem{Génération de datasets synthétiques par fine-tuning inverse itératif de LLMs $\cdot$ \textbf{Python}, \textbf{OpenAI API}, \textbf{Transformers}, \textbf{Pandas} $\to$ \href{https://huggingface.co/blog/zino36/tags-generation-dataset-article}{Article Hugging Face}}
    }
% @endbullets
  }
% @endcvblock

  \mainSection{Formation}
  \cvEntry{Université Paris Cité}
    {Master MLSD -- Machine Learning pour la Science des Données (M1 + M2)}
    {2026 -- 2028}{%
% @bullets id=universite-paris-cite
    \cvBullets{
      \cvItem{\textbf{Alternance sur 24 mois} (M1 + M2, même entreprise) $\cdot$ 3 j université / 2 j entreprise, puis \textbf{temps plein en entreprise dès avril}}
      \cvItem{Au programme \textbf{Deep Learning} $\cdot$ \textbf{IA générative} $\cdot$ \textbf{Text mining \& NLP} $\cdot$ \textbf{Reinforcement Learning} $\cdot$ \textbf{Graph learning} $\cdot$ \textbf{Cloud} $\cdot$ \textbf{Conteneurisation} $\cdot$ \textbf{Big Data} $\cdot$ \textbf{Recommender Systems} $\cdot$ \textbf{Programmation distribuée} $\cdot$ \textbf{Packaging} $\cdot$ \textbf{Projet pluridisciplinaire} (99~h, Centre Borelli)}
    }
% @endbullets
  }

  \cvEntry{Université d'Avignon -- CERI}
    {Licence Informatique, Parcours Intelligence Artificielle}
    {2024 -- 2026}{%
% @bullets id=universite-d-avignon-ceri
    \cvBullets{
      \cvItem{\textbf{Major de promotion} L2 \& L3 (\textbf{1\textsuperscript{ère}/126}) $\cdot$ moyenne \textbf{> 15/20}}
    }
% @endbullets
  }

  \cvEntry{ENSIA -- École Nationale Supérieure d'IA, Alger}
    {\textbf{Classe préparatoire} intégrée au cycle ingénieur -- 120 ECTS, cursus entièrement en anglais}
    {2022 -- 2024}{%
% @bullets id=ensia-ecole-nationale-supe
    \cvBullets{
      \cvItem{Programme sélectif à recrutement national $\cdot$ seuil d'admission 2022 : \textbf{18,29/20} $\cdot$ Mention \textbf{Très Bien}}
    }
% @endbullets
  }

  \cvOneLine{Baccalauréat Sciences Expérimentales {\normalfont\fontsize{9pt}{10pt}\selectfont
    -- Moyenne \textbf{17,82/20} $\cdot$ Mathématiques \textbf{19,5} $\cdot$ Physique \textbf{19}}}{2022}

  \mainSection{Langues, Certifications \& Centres d'intérêt}
  {\fontsize{9pt}{13pt}\selectfont
    \textbf{Anglais} C2 (LanguageCert, certification Ofqual, 02/2026) $\cdot$
    \textbf{Français} bilingue $\cdot$ \textbf{Arabe} maternel\par\vspace{3pt}
    \textbf{AYLP} -- Leadership \& Tech, Nevada (USA) $\cdot$
    \textbf{Hackathon IA Avignon 2025} -- Tech Lead\par\vspace{3pt}
    \textbf{Engagement} Google DSC -- Responsable IT $\cdot$ Mentorat \& ateliers tech $\cdot$
    \textbf{Centres d'intérêt} Piano $\cdot$ Peinture $\cdot$ Tennis \& sport\par
  }

\end{minipage}%
\end{lrbox}%
% Hauteur réelle du contenu, pour que cv_builder la compare à la place disponible.
\setlength{\cvMainHeight}{\dimexpr\ht\cvMainBox+\dp\cvMainBox\relax}%
\typeout{CVFIT content=\the\cvMainHeight\space available=\the\cvMainAvail}%
\usebox{\cvMainBox}%

\end{document}
